\begingroup\setcounter{secnumdepth}{2}

Anomaly detection is a foundational machine learning capability deployed
across critical domains including healthcare monitoring, financial fraud
detection, industrial process supervision, and network intrusion detection.
In each of these domains, the data being analysed is highly sensitive:
patient physiological signals, financial transaction records, process sensor
readings, and network traffic all carry significant privacy requirements that
prohibit direct exposure to third-party cloud servers. Yet conventional anomaly
detection systems require the detection server to have unrestricted access to
raw plaintext data in order to compute anomaly scores, creating a fundamental
and unresolved conflict between anomaly detection and data privacy.

Homomorphic Encryption (HE) offers a cryptographically principled resolution
to this tension. The Homomorphic Encryption scheme permits approximate
arithmetic computation --- addition, multiplication, and their compositions ---
directly on ciphertexts, such that the result, when decrypted by the data
owner, exactly matches the result of the same computation on the original
plaintexts. This property can enable computing anomaly scores on
encrypted data without ever exposing the underlying plaintext values. The
data owner can retain the secret decryption key and receive only the final
anomaly classification decisions.

This thesis presents the complete design, implementation, and network
deployment of a \textbf{Privacy-Preserving Anomaly Detection Framework} built
on the CKKS scheme of Homomorphic Encryption implemented in Microsoft SEAL. The framework addresses the
full pipeline from raw data ingestion to encrypted anomaly scoring and
classification, encompassing preprocessing, CKKS-compatible encoding,
homomorphic computation, network transport, and real-time monitoring.

The central technical challenge in building such a system is that standard
anomaly detection algorithms require division at critical points --- computing
the feature mean $\mu_j = \frac{1}{N}\sum_{i=1}^N x_{ij}$ and the
variance-normalized anomaly score $s(\mathbf{x}) = \sum_j
\frac{(x_j-\mu_j)^2}{\sigma_j^2}$ --- while the CKKS scheme natively supports
only addition and multiplication. This thesis proposes and evaluates three
approaches to encrypted division: Newton-Raphson iterative inversion, Chebyshev
polynomial approximation, and a novel multi-round interactive protocol. The
multi-round protocol, which routes the encrypted denominator through the
key-holding data owner for exact plaintext inversion and re-encryption, is
shown to be the only viable approach, achieving zero approximation error at
zero additional HE depth cost.

Two unsupervised anomaly detection algorithms are implemented entirely under
encryption using the CKKS scheme:

\begin{itemize}
  \item \textbf{HE-SAD (Homomorphic Statistical Anomaly Detection):}
  A complete CKKS implementation of variance-normalized Mahalanobis distance
  scoring, operating through a 9-round interactive protocol. HE-SAD achieves
  AUC $= 0.9533$ on the SWaT Secure Water Treatment benchmark dataset,
  matching the plaintext baseline AUC of $0.929$ and demonstrating that no
  fundamental accuracy penalty results from privacy-preserving encryption.
  The implementation consumes only 2 of 5 available HE multiplication levels
  per test row, leaving substantial headroom for future extensions.

  \item \textbf{HE-PCA (Homomorphic Principal Component Analysis):}
  A complete CKKS implementation of PCA reconstruction error scoring with
  automatic optimal component selection based on a configurable explained
  variance threshold. HE-PCA operates through an 8-round protocol and
  consumes 3 of 5 available HE levels per row. A slight AUC reduction
  compared to the plaintext PCA baseline is observed and attributed to CKKS
  approximation noise accumulating across the $k$ component projection
  operations, a characteristic of leveled HE for matrix computations.
\end{itemize}

The framework is deployed as a network application using FastAPI for the
detection server and Flask with SocketIO for a live browser-based monitoring
dashboard. Ciphertexts are transferred in batches of 100 per HTTP request
to avoid socket buffer overflow. Score computation runs asynchronously in a
background thread. The dashboard embeds full client logic, requiring only two
terminal processes to operate the complete system.

This work is grounded in a prior systematic benchmarking study of three
widely-used FHE libraries --- Microsoft SEAL, IBM HElib, and Pyfhel ---
across four experimental categories covering arithmetic operations, rotation
operations, maximum multiplication depth, and noise budget evolution. The
benchmarking study provides the empirical foundation for library and parameter
selection, identifying SEAL as the most performant library for real-valued
anomaly detection workloads with the deepest available multiplication budget
(depth 32 at polynomial degree 32768).

The framework is evaluated on two publicly available benchmark datasets: SWaT
(51 features, collected from a water treatment testbed) and BATADAL (43
features, from a water distribution network simulation). Both datasets represent
realistic multivariate time-series data with heterogeneous continuous and binary
features, serving as representative benchmarks for privacy-preserving anomaly
detection research.

\endgroup
